% 02-adding-problem.tex ;  The adding problem, expanded with T=100 and d_h=12
%
% Finding: adding problem is NOT a good vanishing-gradient diagnostic;
% even at T=100, d_h=12, gradient ratio stays near 1. The chapter
% uses this to motivate why the copy task is the better test.

\section{Bài toán cộng}
\label{sec:ch02-adding}

\index{adding problem|(}%

\subsection{Định nghĩa}

Cho một chuỗi $T$ bước thời gian. Tại mỗi bước $t$, mạng nhận hai số:
\begin{itemize}
  \item Một giá trị ngẫu nhiên $v_t \sim \mathrm{Uniform}(-1, 1)$.
  \item Một cờ đánh dấu $m_t \in \{0, 1\}$. Cờ bằng $1$ tại \textbf{đúng hai}
    vị trí ngẫu nhiên trong chuỗi; bằng $0$ ở mọi vị trí còn lại.
\end{itemize}

Nhiệm vụ: tại bước cuối cùng $T$, mạng phải xuất ra tổng của hai giá trị tại
hai vị trí được đánh dấu.

\begin{example}
Với $T=10$, hai vị trí được đánh dấu là $t=3$ và $t=7$:
\[
\begin{array}{c|cccccccccc}
t      & 1 & 2 & 3    & 4 & 5 & 6 & 7    & 8 & 9 & 10 \\
\hline
v_t    & 0.5 & -0.3 & \mathbf{0.8} & 0.1 & -0.4 & 0.6 & \mathbf{0.2} & -0.7 & 0.9 & -0.1 \\
m_t    & 0 & 0 & \mathbf{1} & 0 & 0 & 0 & \mathbf{1} & 0 & 0 & 0
\end{array}
\]
Target: $0.8 + 0.2 = 1.0$.
\end{example}

\subsection{Kết quả}

Tôi huấn luyện RNN với $d_h = 24$, learning rate $0.01$, khởi tạo trọng số
$\mathcal{N}(0, 0.1)$, hàm mất mát MSE trên output cuối cùng.

\begin{measured}{Adding problem, $d_h = 24$}
\small
\begin{tabular}{@{}lrrrrr@{}}
\toprule
$T$ & Mẫu & Loss trước & Loss sau &
$\|\partial L / \partial \mat{W}_{hh}\|$ tại $t=1$ & tại $t=T$ \\
\midrule
10  & 5000  & 2.86 & 0.67 & 0.387 & 0.228 \\
20  & 5000  & 1.34 & 0.65 & 0.206 & 0.206 \\
50  & 10000 & 0.002 & 0.66 & 0.235 & 0.196 \\
100 & 20000 & 0.09 & 0.65 & 0.045 & 0.046 \\
\bottomrule
\end{tabular}
\normalsize
\end{measured}

Tôi mong đợi gradient sẽ suy giảm rõ rệt khi $T$ tăng. Nó đã không. Ngay cả
ở $T = 100$, tỉ lệ norm gradient giữa $t=1$ và $t=T$ vẫn gần bằng 1. Loss
giảm từ 2.86 xuống 0.67 ở $T=10$ và giữ nguyên ở mức đó cho đến $T=100$.

Để loại trừ khả năng $d_h = 24$ che giấu vanishing, tôi thử lại với $d_h
= 12$ và $T$ lên đến 100:

\begin{measured}{Adding problem, $d_h = 12$}
\small
\begin{tabular}{@{}lrrrrr@{}}
\toprule
$T$ & Mẫu & Loss trước & Loss sau &
$\|\partial L / \partial \mat{W}_{hh}\|$ tại $t=1$ & tại $t=T$ \\
\midrule
50  & 10000 & 1.28 & 0.69 & 0.012 & 0.012 \\
100 & 20000 & 0.03 & 0.67 & 0.045 & 0.046 \\
\bottomrule
\end{tabular}
\normalsize
\end{measured}

Vẫn không có vanishing. Với $d_h = 12$ và $T = 100$, tỉ lệ gradient là
1.01; về cơ bản là phẳng.

Đây không phải là thất bại của thí nghiệm. Đây là một kết quả quan trọng:
\textbf{bài toán cộng không phải là bài kiểm tra tốt cho vanishing gradient}.
Lý do: chỉ output cuối cùng có target, nên backward signal chỉ xuất phát từ
$t = T$, và nó chỉ cần đi qua $T$ bước \emph{một lần}. Hàm cần học; tổng hai
số; quá đơn giản để tạo ra áp lực lên gradient.

Bài toán cộng kiểm tra xem RNN có học được không. Bài toán tiếp theo kiểm tra
xem gradient có sống sót không.

\index{adding problem|)}%
